\documentclass[11pt,a4paper]{article}

\usepackage[margin=1in]{geometry}
\usepackage{amsmath,amssymb}
\usepackage{listings}
\usepackage{xcolor}
\usepackage{booktabs}
\usepackage{enumitem}
\usepackage{parskip}
\usepackage{fancyhdr}

\pagestyle{fancy}
\fancyhf{}
\rhead{CS335 Project Milestone Report}
\lhead{Relational Verification of Hyperproperties}
\rfoot{\thepage}

\definecolor{codebg}{rgb}{0.95,0.95,0.95}
\definecolor{codegreen}{rgb}{0,0.5,0}
\definecolor{codepurple}{rgb}{0.5,0,0.5}
\definecolor{codeblue}{rgb}{0,0,0.7}

\lstdefinestyle{code}{
  backgroundcolor=\color{codebg},
  basicstyle=\ttfamily\small,
  keywordstyle=\color{codeblue}\bfseries,
  commentstyle=\color{codegreen}\itshape,
  frame=single,
  rulecolor=\color{gray!40},
  breaklines=true,
  tabsize=4,
  showstringspaces=false,
}

\title{
  \Large\textbf{CS335 Project Milestone Report}\\[0.4em]
  \large Relational Verification of Hyperproperties in Chiron
}
\author{}
\date{\today}

\begin{document}
\maketitle
\thispagestyle{fancy}

% ============================================================
\section{Overview}
% ============================================================

Chiron provides several single-trace analyses but cannot reason about
\textbf{hyperproperties} — properties relating multiple executions of the same
program. We add a relational verification pass that checks
\textbf{non-interference}: two executions agreeing on all public (low) inputs
must produce identical public outputs, regardless of their secret (high) inputs.

The core technique is \textbf{self-composition}: encode both executions into one
Z3 query, constrain low inputs to agree, and ask whether low outputs can differ.
\textsc{unsat} proves non-interference; \textsc{sat} yields a concrete
counterexample.

Formally, for a program $P$ with state partitioned into low/high:
\[
  \forall \sigma_1, \sigma_2.\;
  \sigma_1|_{L} = \sigma_2|_{L}
  \;\Rightarrow\;
  \llbracket P \rrbracket(\sigma_1)|_{L} = \llbracket P \rrbracket(\sigma_2)|_{L}
\]

We check the negation: find $\sigma_1, \sigma_2$ with same low inputs but
different low outputs. \textsc{unsat} means no such pair exists.

% ============================================================
\section{Implementation}
% ============================================================

\subsection{Files Changed}

\begin{itemize}[noitemsep]
  \item \textbf{\texttt{ChironCore/relationalVerifier.py}} (new): Core verifier
  implementing both tiers.
  \item \textbf{\texttt{ChironCore/chiron.py}} (modified): Added \texttt{-rv}
  (Tier 1), \texttt{-rvl} (Tier 2), \texttt{--low\_in}, and \texttt{--low\_out}
  CLI flags.
  \item \textbf{\texttt{ChironCore/example/}} (new files): 12 test programs
  covering safe and leaky cases for both tiers.
\end{itemize}

\subsection{Key Design: Reusing Chiron's Symbolic Encoding}

Chiron's existing \texttt{handleAssignment} in \texttt{sExecutionInterface.py}
symbolically evaluates an assignment by building a Z3 expression:

\begin{lstlisting}[style=code, language=Python]
def handleAssignment(z3Vars, stmt):
    lhs = str(stmt.lvar).replace(":", "")
    rhs = str(stmt.rexpr).strip().replace(":", "z3Vars.")
    setAttr(z3Vars, lhs, convertExp(z3Vars, rhs))
\end{lstlisting}

It works on any context object passed as \texttt{z3Vars}. We create \emph{two}
such contexts with differently-suffixed Z3 variable names (\texttt{pub\_1},
\texttt{pub\_2}, etc.) and run the same encoding logic on both, giving us two
independent symbolic traces in one Z3 namespace.

% ============================================================
\section{Tier 1: Straight-Line Programs}
% ============================================================

For programs with no loops or conditionals, the entire verification is a single
Z3 \texttt{check()} call:

\begin{enumerate}[noitemsep]
  \item Create \texttt{ctx1} (suffix \texttt{\_1}) and \texttt{ctx2} (suffix
        \texttt{\_2}) with free Z3 integer variables for each program variable.
  \item Walk the IR with \texttt{encode\_stmts} for both contexts. Each
        \texttt{ctx.var} ends up holding the final Z3 formula for \texttt{:var}.
  \item Add: $\texttt{pub\_1} = \texttt{pub\_2}$ for all low inputs.
  \item Add: $\bigvee_{v \in L_{out}} v_1 \neq v_2$ (can outputs differ?).
  \item \texttt{s.check()} — \textsc{unsat} proves safety; \textsc{sat} gives a
        counterexample.
\end{enumerate}

\subsection*{Examples Tested}

\begin{table}[h!]
\centering\small
\begin{tabular}{llcp{5.5cm}}
\toprule
\textbf{File} & \textbf{Program} & \textbf{Result} & \textbf{Why} \\
\midrule
\texttt{ni\_safe.tl}  & \texttt{:out = :pub + 1}                              & \textsc{unsat} & Secret unused \\
\texttt{ni\_safe2.tl} & \texttt{:a=:s+:p; :b=:a-:s; :out=:b}                 & \textsc{unsat} & Secret cancels out \\
\texttt{ni\_safe3.tl} & \texttt{:san=:s-:s; :out=:pub+:san}                   & \textsc{unsat} & Sanitised to 0 \\
\texttt{ni\_safe4.tl} & \texttt{:x=:s*5; :y=:p+:x; :z=:y-:x; :out=:z}       & \textsc{unsat} & 3-step cancellation \\
\midrule
\texttt{ni\_leak.tl}  & \texttt{:out = :pub + :secret}                        & \textsc{sat}   & Direct flow \\
\texttt{ni\_leak2.tl} & \texttt{:t=:s*2; :out=:pub+:t}                        & \textsc{sat}   & Via temp variable \\
\texttt{ni\_leak3.tl} & \texttt{:a=:p+1; :b=:s-:p; :out=:a+:b}               & \textsc{sat}   & Hidden: \texttt{out=s+1} \\
\texttt{ni\_leak4.tl} & \texttt{:a=:s+3; :b=:a*2; :c=:b-6; :out=:p+:c}      & \textsc{sat}   & 3-level indirection \\
\bottomrule
\end{tabular}
\end{table}

% ============================================================
\section{Tier 2: Programs with a Single Loop}
% ============================================================

For \texttt{repeat :n [body]}, the iteration count $n$ is a symbolic unknown.
Unrolling is not possible. Instead we use a \textbf{relational loop invariant}
$\mathit{INV}(\sigma_1, \sigma_2)$ — a predicate that holds at the loop head at
every iteration — and prove it inductively via three Z3 queries.

\subsection*{Default Invariant}

We automatically derive the invariant: all low variables and the loop counter
are equal across both traces at every iteration.
\[
  \mathit{INV}(\sigma_1, \sigma_2) =
  \bigwedge_{v \in \mathit{Low}} (\sigma_1(v) = \sigma_2(v))
  \;\wedge\;
  (\mathit{counter}_1 = \mathit{counter}_2)
\]

\subsection*{Three Verification Checks}

\paragraph{Check 1 — Initialisation.} INV holds at loop entry (given equal low
inputs and the symbolically encoded pre-loop code).

\paragraph{Check 2 — Preservation.} With \emph{fresh} unconstrained Z3
variables for the loop-head state (constrained only by INV and the loop
condition), one body execution does not break INV.
\[
  \mathit{UNSAT}\bigl(
    \mathit{INV}(\sigma_1,\sigma_2)
    \wedge \mathit{cond}(\sigma_1)
    \wedge \mathit{cond}(\sigma_2)
    \wedge \neg\,\mathit{INV}(\mathit{body}(\sigma_1), \mathit{body}(\sigma_2))
  \bigr)
\]

\paragraph{Check 3 — Consequence.} At loop exit, INV holds and the loop
condition is false. After encoding any post-loop code, low outputs cannot
differ.

All three \textsc{unsat} $\Rightarrow$ non-interference holds for any $n$.
Each failing check emits a concrete counterexample from Z3's model.

\subsection*{Examples Tested}

\begin{table}[h!]
\centering\small
\begin{tabular}{lp{5cm}cc}
\toprule
\textbf{File} & \textbf{Program} & \textbf{Check 2} & \textbf{Check 3} \\
\midrule
\texttt{loop\_safe.tl}         & \texttt{:out=0; repeat :n [:out=:out+:pub]}                   & pass & pass \\
\texttt{loop\_leak.tl}         & \texttt{:out=0; repeat :n [:out=:out+:secret]}                & fail & ---  \\
\texttt{loop\_postloop\_safe.tl} & safe loop $+$ \texttt{:out=:out+:pub} after                & pass & pass \\
\texttt{loop\_postloop\_leak.tl} & safe loop $+$ \texttt{:out=:out+:secret} after             & pass & fail \\
\bottomrule
\end{tabular}
\end{table}

\texttt{loop\_postloop\_leak.tl} highlights why post-loop code must be encoded
in Check 3: the loop body is safe, but the post-loop assignment leaks the
secret, caught only by the consequence step.

% ============================================================
\section{Future Work}
% ============================================================

\paragraph{Non-trivial invariants.}
The default invariant fails for programs where secrets appear in intermediate
variables that contribute to public outputs, even if they ultimately cancel. For
example:
\begin{lstlisting}[style=code]
:temp = :secret + :pub
repeat :n [ :out = :out + :temp - :secret ]
\end{lstlisting}
This is safe (\texttt{out = (n+1)*pub}) but \texttt{temp} depends on
\texttt{:secret}, so \texttt{temp\_1 $\neq$ temp\_2}. The correct invariant
needs the cross-trace relation $\mathit{temp}_1 - \mathit{secret}_1 =
\mathit{temp}_2 - \mathit{secret}_2$, which cannot be inferred automatically.
Tier 3 will support user-provided invariant formulas for such cases.

\paragraph{Conditionals and implicit flows.}
Programs with \texttt{if :secret > :pub [...]} inside loops exhibit implicit
information flows: the secret controls which branch is taken, leaking
information through the output even without a direct assignment. Handling this
requires encoding conditionals as Z3 \texttt{If} expressions and reasoning
about trace pairs that diverge at branches.

\paragraph{Multiple and nested loops.}
The current implementation handles one loop per program. Multiple sequential
loops will be handled by chaining the three-check procedure, using the verified
exit state of one loop as the entry state of the next.

\end{document}